% Generated draft: requires mathematical and source review.
\documentclass[11pt]{article}
\usepackage[margin=25mm]{geometry}
\usepackage{fontspec}
\setmainfont{Latin Modern Roman}
\usepackage{amsmath,amssymb,mathtools,unicode-math}
\setmathfont{Latin Modern Math}
\usepackage{enumitem,needspace,xurl,hyperref}
\setlength{\parindent}{0pt}
\setlength{\parskip}{5pt}
\setlength{\emergencystretch}{3em}
\newcommand{\N}{\mathbb{N}}
\newcommand{\Z}{\mathbb{Z}}
\newcommand{\Q}{\mathbb{Q}}
\newcommand{\R}{\mathbb{R}}
\newcommand{\C}{\mathbb{C}}
\newcommand{\F}{\mathbb{F}}
\newcommand{\A}{\mathbb{A}}
\newcommand{\PP}{\mathbb{P}}
\newcommand{\E}{\mathbb{E}}
\newcommand{\eps}{\varepsilon}
\newcommand{\dd}{\,\mathrm{d}}
\begin{document}
\section*{UnsolvedMath ID: 30004735}
\textbf{Source problem code:} catalogue entry
\subsection{Definitions and mathematical statement}
Let \(\Omega\) be a billiard domain and \(x\in\partial\Omega\). The supplied record does not specify the ambient geometry, dimension, boundary regularity, or convexity assumptions on \(\Omega\). For a positive integer \(q\), write formally
\[
\mathcal L_q(\Omega,x)=\left\{\gamma:\ \gamma\text{ is a billiard geodesic loop based at }x,\ \operatorname{wind}(\gamma)=1,\ N_{\mathrm{refl}}(\gamma)=q\right\}.
\]
Here a loop based at \(x\) begins and ends at \(x\); \(N_{\mathrm{refl}}\) denotes the number of reflections under the intended counting convention, and \(\operatorname{wind}\) denotes the intended winding number. These last two conventions are not defined in the supplied record. In the usual Euclidean interpretation, a billiard path consists of straight segments joined at boundary contacts satisfying the equal-angle reflection law wherever that law is defined. This interpretation alone does not resolve the missing hypotheses or conventions.

\textbf{Question, conditional on specifying the admissible class and conventions.} For every domain \(\Omega\), boundary point \(x\), and reflection count \(q\) in the intended ranges, does there exist exactly one such loop, that is,
\[
\#\mathcal L_q(\Omega,x)=1?
\]
If its unique member is denoted by \(\gamma_{q,x}\), the proposed loop-length function is
\[
\psi_q(x)=\operatorname{Length}(\gamma_{q,x}),\qquad x\in\partial\Omega.
\]
The assertion is an existence-and-uniqueness question, not a theorem asserted here. The source leaves unspecified the allowable values of \(q\), the meaning of winding number, whether the base point counts as a reflection, whether the reflection law is imposed at the closing contact, and when two loops are identified, including under reparametrization or reversal. Consequently, the displayed formulation is a conditional formalization of the supplied target; it is not a fully specified conjecture for arbitrary billiard domains. Uniqueness of a loop would make its length unambiguous, but a common length for several loops would be a weaker property.
\subsection{Short English statement}
The problem asks whether each boundary point of a billiard domain supports exactly one loop with winding number \(1\) and \(q\) reflections, thereby defining its length \(\psi_q(x)\). A precise conjecture requires additional domain hypotheses and conventions for winding number, reflection counting, and equality of loops. The supplied registry labels the problem partially solved; that status is not independently verified here.
\subsection{Sources}
\begin{itemize}
\item Supplied local UnsolvedMath registry record, permanent ID \(30004735\), source problem code \texttt{catalogue entry}, title: \emph{Uniqueness of Fixed-Bounce Geodesic Loops in Billiards}. The record belongs to the Oberwolfach Reports Open Problems collection. No individual proposer is supplied; the proposed year is recorded as 2021.
\item Supplied report citation: \emph{Dynamische Systeme (2022)}, Oberwolfach Reports, DOI 10.4171/owr/2021/33, pp.\ 1735--1803. \url{https://doi.org/10.4171/owr/2021/33}. The supplied background describes these reports as documenting mathematical workshops and open-problem discussions. No precise page or passage for this problem is supplied.
\item Supplied literature reference, labeled ``analytic symmetric domains'': \url{https://annals.math.princeton.edu/2009/170-1/p06}.
\item Supplied literature reference, labeled ``marked-length determination'': \url{https://doi.org/10.1007/s00222-023-01191-8}.
\item Supplied literature reference, labeled ``convex spectral invariants'': \url{https://doi.org/10.4310/jdg/1214437138}.
\item Literature assessment checked at \texttt{2026-08-22}: this is supplied dated curation, not independently verified. It labels the problem partially solved and states that analytic symmetric and selected chaotic billiards admit marked-length or spectral determination results, while general well-definedness of the loop-length data and spectral detection of convexity remain unresolved. The background explicitly calls this dated literature triage derived from the supplied OWR curation and requests independent verification of time-sensitive status and cited claims. Spectral detection of convexity is background context, not an additional assertion in the displayed loop-uniqueness question.
\item The registry labels the statement \texttt{corrected\_verified} and reports a manual semantic, typographic, and LaTeX review claiming self-containedness and preservation of the mathematical target. This is supplied review metadata, not independent verification of the formulation or its mathematical correctness.
\end{itemize}
\end{document}
